The seats-votes curve $S(v)$---the function mapping statewide Democratic two-party vote share
to expected Democratic seat share under uniform swing---is the sufficient statistic for partisan
fairness from which both the Efficiency Gap and Partisan Bias are analytically recoverable.
Efficiency Gap equals the signed area between the $S(v)$ curve and the diagonal $S(v) = v$;
Partisan Bias equals $S(0.50) - 0.50$. Responsiveness, defined as $R = dS/dv|_{v=0.50}$,
is the additional summary statistic measuring how many seat-share points change per vote-share
point at the electoral pivot. A competitively ideal map has $R \approx 2.0$; suppressed
competition ($R < 1.5$) indicates a high proportion of safe seats insulating incumbents
from vote swings.

We find that \textsc{bisect} algorithmic maps have $R \approx 2.0$ for NC ($k=14$) and
WI ($k=8$)---the competitive ideal---because compactness-driven bisection produces a roughly
uniform distribution of district vote shares without systematic concentration in safe-seat
territory. Enacted NC 2022 congressional maps have $R \approx 1.3$, reflecting the high
proportion of safe Republican seats created by packing Democratic voters and the resulting
insulation of Republican incumbents from vote swings.

The seats-votes curve is the preferred single courtroom exhibit for partisan fairness evidence:
it displays both bias ($S(0.50)$ vs.\ 0.50) and responsiveness ($dS/dv$) in a single graph,
it communicates intuitively to judicial audiences without requiring them to understand
wasted-vote algebra, and it has been admitted as expert evidence in NC \textit{Harper v.\ Hall}
and OH redistricting proceedings. Comparing the bisect $S(v)$ curve (passing through $(0.50, 0.50)$
with slope $\approx 2.0$) to the enacted curve (passing below $(0.50, 0.50)$ with slope $\approx 1.3$)
provides the most complete single visual exhibit of partisan asymmetry available.
